\begin{abstract}
Emergency departments (EDs) differ in case mix, load and layout, yet streaming policies, that is, the order in which patients are served and the rule that sends them to treatment areas, are often copied between EDs. This paper asks which policy fits which type of ED, how much is lost when one policy is used everywhere, how accurate predictions must be for prediction-based policies to pay off, whether similarity-based scoring rules apply the policies they are meant to represent, and how much a scoring rule tuned to an ED can gain. A discrete-event simulation of \NEdTypes{} synthetic ED types (three acuity mixes, three load levels, three layouts) compares seven queue rules and three routing rules under common random numbers. The common default, triage-level order with a triage-based fast track, had a mean regret of \DefaultRatpMeanRegret\% in an acuity-weighted breach rate. A single well-chosen policy came close to the best for any one measure, but which policy depended on the measure. Prediction-based policies gained mainly through rule structure, not prediction accuracy. We state conditions under which a similarity rule equals a classical rule; when they fail, average outcomes can match while decisions differ at about half of all dispatches. Similarity rules tuned to each ED type by a cross-entropy search and evaluated on separate test replications improved on the best classical policy by \OptBalancedGainBestMean\% on average in a balanced objective and by up to \OptBalancedGainBestMax\% in split-flow EDs. The conclusions held under misspecified inputs and in \NOosTypes{} further ED types.
\end{abstract}

\begin{IEEEkeywords}
Emergency department, patient streaming, prioritization, discrete-event simulation, policy fit, regret, decision support, healthcare operations management.
\end{IEEEkeywords}

\noindent\textbf{Managerial relevance statement}---ED and health system managers decide how waiting patients are prioritized and whether to run a fast track or split flow, and regional programmes often standardize these choices across sites. This study shows that the common practice of strict triage order with a triage-based fast track can fall well short of the best available policy, and that no single policy is best for every objective. Managers should first state which objective matters most (timely access for urgent patients, overall length of stay, or the longest waits of low-acuity patients) and then choose the policy that fits their ED's case mix and layout; a decision guide is provided. When streaming is introduced, the routing rule must match the split of staff and spaces between areas, or one area can be overloaded. Prediction tools add value mainly through how they change routing. A scoring rule can be tuned to an ED by simulation; tuning should use a balanced objective, because optimizing access alone delayed reassessments and lengthened stays. Weights of a scoring rule should be checked against the intended policy at the level of individual decisions, because similar averages can hide different choices for patients.
